\scanpage{199}

% B01 / text / top, section 3
\mspage{102}3.\quad Je vais reprendre la notion \inserted{de} graphe et de graphe triangulaire,
dans une optique un peu différente. Une structure
de graphe strict peut être interprétée comme la donnée
d'un \deleted{\unsure{gr}} ens. de sommets $S$ et d'une partie
\uline{symétrique} de $S\times S\smallsetminus S$

% B02 / display / upper, (1) and NB to the left
\[
 R_0\subset S\times S\unsure{\smallsetminus|\vphantom{S}},\quad R_0\cap\mathrm{diag}=\varnothing,\quad{}^sR=R, \tag*{(1)}
\]
% slanted
\margin{\uline{NB} si on n'exige pas $R_0\cap\mathrm{diag}=\varnothing$, on trouve une notion un peu plus générale, avec $A_0\subset\unsure{P|\mathcal P|\mathfrak P}_{\mathrm{card}\in\{1,2\}}(S)$, qui correspond à celle d'un graphe pouvant avoir des boucles fermées (une au plus par sommet)}

% B03 / text / upper, before (2)
l'étoile d'un $s\in S$ s'écrivant alors comme

% B04 / display / upper, (2), notation question on right
\[
 \mathrm{Et}_{A_0}(s)=R_0(s)\qquad \text{on note aussi }\mathrm{Et}_{R_0}(s)\ ? \tag*{(2)}
\]

% B05 / text / upper, before (3)
\unsure{puis|que|pour} l'on a

% B06 / display / upper middle, (3)
\[
 A_0=\{\{\unsure{s,t|s;t}\}\in P_2(S)\mid(s,t)\in R_0\}\ . \tag*{(3)}
\]

% B07 / text / middle, restart before (4)
\deleted{Dans le cas où} \inserted{Considérons} l'application

% B08 / diagram / middle, (4)
\diagram{(4)\qquad\begin{tikzcd}[ampersand replacement=\&,row sep=0pt]
 S\arrow[r,mapsto]\&P(S)\\ s\arrow[r,mapsto]\&R_0(s)=\mathrm{Et}_{A_0}(s)
\end{tikzcd}}

% B09 / text / middle, introduction of the image
\deleted{est injective,} désignons \deleted{en tous cas} par $S^*$ son \deleted{\unsure{image}}\inserted{image},
(notée $\unsure{\mathcal E|\mathscr E|\Sigma}$ précédemment

% B10 / display / middle, (5)
\[
 S^*=\{E\in P(S)\mid\exists\,s\in S,\ \text{t.q. }E=R_0(s)\}, \tag*{(5)}
\]

% B11 / text / middle, before (6)
on note que l'inclusion identique

% B12 / display / middle, (6), inclusion and assignment
\[
 S^*\lhook\joinrel\longrightarrow P(S)\qquad s^*\longmapsto E_{s^*} \tag*{(6)}
\]

% B13 / text / lower middle, interpretation
s'interprète aussi comme définissant une relation

% B14 / display / lower middle, (7)
\[
 \deleted{\unsure{\rho|P}}\unsure{R_*|R_{\star}|\vec{R}_*}\subset S\times S^* \tag*{(7)}
\]

% B15 / text / lower, before (8)
liée à (6) par

% B16 / display / lower, (8)
\[
 \forall\,s^*\in S^*,\quad\text{on a}\quad{}^sR_*(s^*)=E_{s^*}\ . \tag*{(8)}
\]

% B17 / text / lower, bracketed special case
[\deleted{d'où} Nous nous intéressons surtout au cas où

% B18 / display / lower, crossed (9)
\[
 \deleted{\forall\,s\in S,\ \text{on a }R_0(s)} \tag*{\deleted{(9)}}
\]

% B19 / text / lower, (4) injective
(4) est injective (\deleted{\unsure{tout}}\inserted{\unsure{une|un} ($A_0$-)} étoile détermine son centre),
et où on \deleted{peut écrire}\inserted{trouve} alors, pour tout $s\in S^*$, une
application bijection canonique

% B20 / diagram / foot, (9)
\diagram{(9)\qquad\begin{tikzcd}[ampersand replacement=\&,row sep=small]
 \underbrace{R_*(s)}_{\substack{\text{ens des étoiles}\\ E\ \text{qui contiennent }s}}\&\mathrm{Et}_{A_0}(s)\arrow[l,"\sim"']\\
 \mathrm{Et}_{A_0}(t)\&t\arrow[l,mapsto]
\end{tikzcd}}

% B21 / text / foot, last parenthesis
(en général, on a seulement une application surjective)]

% archive sheet number: 198
